% ============================================================================
% Chapter 23 --- Wave 5: TAC pilot
% ----------------------------------------------------------------------------
% Purpose : exercise the full stack on a contained scenario. Absorbs the
%           original A.3 (minimum-viable pilot) content. First wave that
%           exercises the architecture end-to-end.
% Page target: ~6 pages.
% Dependencies: Waves 1, 2, 3, 4 closed (D17.3); reference layers Ch. 6-14.
% Adapts wave template (D18.2): three illustrative pilot scenarios in
% addition to (or in place of) the three starting-position trajectories,
% since the pilot's defining variable is the scenario chosen, not the
% institution's starting position.
% ----------------------------------------------------------------------------
% Decisions registered in _coherence-EN.md:
%   D23.1 Wave 5 adapts the wave template: scenarios in place of trajectories
%   D23.2 three illustrative pilot scenarios published
%   D23.3 pilot lives in an isolated zone, no production impact
%   D23.4 exit criteria reference the union of all conformance subsets
%         exercised by the pilot
% ============================================================================

\chapter{Wave~5 --- TAC pilot}
\label{ch:wave5-pilot}

\begin{caveatbox}[Dependencies and template adaptation]
\noindent Wave~5 applies the wave chapter template established in
\trchap{ch:wave0-foundations}{} with one adaptation: the principal
differentiating variable is the chosen pilot scenario rather than
the institution's starting position, so the chapter develops three
illustrative scenarios (\S\ref{sec:w5-scenarios}) and consolidates
the trajectory discussion (\S\ref{sec:w5-trajectories}). Required
pre-requisites are Waves~1, 2, 3 and 4. Closure is gated by the
union of conformance subsets corresponding to the layers the pilot
exercises.
\end{caveatbox}

Wave~5 is the wave in which the architecture is exercised
end-to-end, in production conditions, on a contained scope. It is
the first wave that demonstrates the integration of the
preceding waves, and the first opportunity for the institution to
encounter the architecture's emergent behaviour --- the
combination effects that no single layer's testing can surface.

%-----------------------------------------------------------------------------
\section{Goal and dependencies}
\label{sec:w5-goal}

The goal of Wave~5 is to operate one or more Trusted Activity
Contexts in production, on a scope contained enough that failure
of the pilot does not affect institutional services, and to
extract from the experience the learnings that subsequent waves
(Endpoint, Audit-and-IR, Workload, Lifecycle) will apply at
larger scale. The pilot is consequently both a technical
demonstration and an organisational learning event.

\paragraph{Pilot, not deployment.} The wave does not promise
generalisation of the chosen scenario across the institution. A
pilot that succeeds operationally has produced what the wave
required; further deployment is the work of subsequent waves and
of the institution's own scaling decisions.

%-----------------------------------------------------------------------------
\section{Pre-requisites}
\label{sec:w5-prereq}

Beyond the closure of Waves~1 to 4, three operational conditions
strengthen Wave~5:

\begin{itemize}[leftmargin=1.5em,nosep]
  \item a participating faculty, research project, or service that
    has accepted to host the pilot, with named accountable
    contacts who will engage with the engineering team
    throughout;
  \item a pilot-specific governance arrangement, typically a small
    steering group composed of the pilot host, the programme
    sponsor, and the engineering lead, meeting at a higher cadence
    than the broader programme governance;
  \item a documented success-and-failure framework agreed at wave
    open, so that the pilot's outcome can be evaluated against
    declared criteria rather than against retrospective
    interpretation.
\end{itemize}

%-----------------------------------------------------------------------------
\section{Coexistence pattern}
\label{sec:w5-coexistence}

Wave~5's coexistence pattern is structurally simpler than the
preceding waves: the pilot operates in an isolated zone of the
network and on isolated tenancies of the compute and storage
layers, without affecting production services. Coexistence is
spatial rather than temporal --- the pilot and the production
services share the institution's infrastructure but do not share
critical operational paths.

\paragraph{Bounded blast radius.} The architectural commitment of
Wave~5 is that an incident within the pilot --- a misconfiguration,
a security event, a performance regression --- does not propagate
to production. The mechanisms that bound the blast radius (zone
isolation, separate identity attribution, separate audit
correlation) are part of the pilot's design.

%-----------------------------------------------------------------------------
\section{Three illustrative pilot scenarios}
\label{sec:w5-scenarios}

The three scenarios below are illustrative; institutions choose
the scenario that aligns with their strategic priorities, with
their willingness to engage the pilot host, and with the
sovereignty grid's current emphasis. The scenarios are not
exclusive --- an institution may run more than one pilot in
sequence or in parallel --- but a single first pilot is generally
preferable to ensure the team's full attention.

\subsection{Scenario A: Researcher-LLM-API}
\label{sec:w5-scenario-llm}

A research group requests scoped, audited access to a
generative-AI cloud endpoint from a dedicated research zone, with
data classification rules and a documented egress policy. The
scenario surfaces most cross-cutting concerns at once: federated
identity, policy decision combining role and posture, network
attachment to a research zone, third-country data egress under
\loc{cross-border-transfer-regime} (in the EU, \gls{gdpr}
Chapter~V, Arts.~44--49, read with the \emph{Schrems~II} doctrine),
audit
aggregation across layers.

\paragraph{Activation sequence.} The researcher authenticates
through the institutional SSO and requests activation of a
\emph{Researcher-LLM-API} role. The identity layer issues a
short-lived assertion. The policy layer evaluates the request
against the institutional rules: role granted, posture verified,
classification of the data the researcher intends to access
permitted under the egress policy. The network layer attaches the
researcher's session to the research zone with the appropriate
egress allowance. The compute layer provisions or attaches the
research workspace. The storage layer binds the project volumes
with the classification labels active. Every step emits events
under the canonical TAC schema.

\paragraph{Why this scenario.} The cross-jurisdiction data flow is
the most institutionally consequential aspect of the scenario. A
pilot that demonstrates a documented egress policy under the
conditions of \loc{cross-border-transfer-regime} (in the EU,
\gls{gdpr} Chapter~V read with the \emph{Schrems~II} doctrine)
produces an artefact that the
institution can show to its governance body, its data-protection
function, and the regulators that may inspect its operations.

\subsection{Scenario B: Teaching cyber-range}
\label{sec:w5-scenario-cyber}

A teaching department requests a time-bounded laboratory
environment for cybersecurity instruction: students receive
isolated networks within a controlled zone, instructors hold
elevated roles for the class duration, and the entire environment
is decommissioned when the term ends.

\paragraph{Why this scenario.} The teaching cyber-range exercises
the lifecycle of a TAC --- activation at term start,
re-activation per class session, decommissioning at term end ---
more rigorously than other scenarios. It also exercises
multi-tenant isolation at a scale (typically tens of student
contexts simultaneously active) that other pilots do not.

\subsection{Scenario C: Administrative compliance workspace}
\label{sec:w5-scenario-compliance}

An administrative service handling regulated data (research-ethics
reviews involving sensitive subject data, compliance reporting
toward regulatory bodies) requests a workspace with stricter
audit, restricted egress, and disclosure-regime guarantees. The
scenario exercises the higher-classification end of the
data-dimension spectrum.

\paragraph{Why this scenario.} The compliance scenario exercises
HYOK arrangements (\trchap{ch:storage}{} D10.6), the strictest
audit-retention obligations, and the integration with existing
institutional compliance processes. It is the scenario that most
directly engages the institution's governance bodies beyond the IT
function.

%-----------------------------------------------------------------------------
\section{Trajectories (consolidated)}
\label{sec:w5-trajectories}

Across the three principal trajectories of the playbook, Wave~5
varies in its integration burden but not in its essential form.

\begin{itemize}[leftmargin=1.5em,nosep]
  \item \textbf{Greenfield.} The pilot is the institution's first
    operational TAC; the integration burden is minimal because
    there is no incumbent to coexist with. The wave is
    correspondingly compressed (indicative duration four to six
    months).
  \item \textbf{Brownfield single-suite-centric.} The pilot operates
    against the institution's existing incumbent identity suite and
    office-productivity estate (e.g.\ Entra-based identity with a
    Microsoft~365 estate), with policy decisions evaluated by the
    Wave~4 OPA deployment. Integration testing focuses on the
    incumbent-IdP-to-policy-plane decision chain.
  \item \textbf{Brownfield hybrid.} The pilot operates against
    the existing Shibboleth federation and the Wave~1 Keycloak
    activation layer. Integration testing focuses on the
    cross-component activation sequence.
\end{itemize}

\paragraph{Indicative duration.} Six to nine months for brownfield
trajectories, with the variability driven by the chosen scenario
and by the institutional change-management cadence.

%-----------------------------------------------------------------------------
\section{Reversibility criteria}
\label{sec:w5-reversibility}

Wave~5 is structurally reversible: the pilot can be stopped at
any time without affecting production services, and the data
generated during the pilot is preserved according to the
institutional classification scheme.

\begin{itemize}[leftmargin=1.5em,nosep]
  \item Pilot stop procedure: documented, exercised in tabletop
    form before pilot start, and invoked without escalation when
    the pilot host requests it.
  \item Data preservation: research data, teaching artefacts, or
    administrative artefacts produced during the pilot are
    exported in documented formats before the pilot
    decommissions.
  \item Learning preservation: the technical and organisational
    learnings from the pilot are documented in a wave-closure
    report whose content is preserved regardless of whether the
    pilot was operationally successful.
\end{itemize}

%-----------------------------------------------------------------------------
\section{Exit criteria}
\label{sec:w5-exit}

Wave~5 closes when the conformance subsets corresponding to the
layers exercised by the pilot pass on the pilot configuration,
when the pilot host has reported its operational experience, and
when the wave-closure report has been reviewed by the governance
body.

\begin{itemize}[leftmargin=1.5em,nosep]
  \item Conformance: the union of the relevant Identity, Policy,
    Network, Compute, Storage and Observability tests passes
    against the pilot configuration. The exact subset depends on
    the chosen scenario.
  \item User experience: the pilot host has produced a structured
    report describing the operational experience --- what worked,
    what did not, what should be revised before any broader
    deployment.
  \item Wave-closure report: technical learnings, organisational
    learnings, and recommendations for subsequent waves are
    consolidated and reviewed by the programme governance.
\end{itemize}

\begin{realworldbox}[Research-cloud pilot precedents]
\noindent Indiana University's Jetstream2, operated under the
United States National Science Foundation's ACCESS programme, is
a publicly documented research-cloud architecture from which
pilot patterns can be drawn. European institutions participating
in the EGI federated cloud and in EUDAT-aligned research-data
projects have produced comparable patterns. The institutional
specifics --- pilot scope, success criteria, durations --- vary
substantially across institutions and should be evaluated against
each institution's own context rather than transposed wholesale.
\end{realworldbox}

%-----------------------------------------------------------------------------
\section{Regulatory anchoring}
\label{sec:w5-regulatory}

\noindent Wave~5 is the first wave that exercises the
architecture end-to-end on a research-class TAC. Its regulatory
profile depends on the pilot scenario chosen. The
Researcher-LLM-API scenario is the principal locus where
\loc{ai-regulation} bites; in the EU, the applicable frame for
commercial-API use is the AI Act's risk-based classification, with
Art.~50 (transparency obligations) applicable when end-users
interact with AI-generated content --- the Art.~2(6)/2(8) research
exemptions cover bespoke, pre-market research tools, not the
consumption of a commercial API.
For all pilot scenarios, \gls{gdpr}~Art.~25 (data protection by
design) and Art.~35 (DPIA where the pilot processes personal
data on a non-trivial scale) frame the documentation duty. The
wave is the first integration test for the
\gls{gdpr}-Art.~32 / NIS2-Art.~21(2)(a) measures introduced by
Waves~1--4. Detailed mapping in
\trchap{ch:regulatory-mapping}{} \S\ref{sec:rm-gdpr},
\S\ref{sec:rm-nis2}, \S\ref{sec:rm-ai-act}.
